% =============================================================================
%  Notas del profesor — Sesión 10: Revisión del Examen Parcial 01
%  Programación Avanzada — Universidad Panamericana
%  Compilar: pdflatex sesion10_revision_examen_profesor.tex
% =============================================================================
\documentclass[12pt, oneside]{article}
\usepackage[profesor]{progavanzada}

\begin{document}

\portada{Sesión 10}{Revisión del Examen Parcial 01}{2026}

\tableofcontents
\newpage

% =============================================================================
\section{Objetivo de la sesión}
% =============================================================================

Regresar los exámenes calificados y revisar cada pregunta en el pizarrón.
No hay contenido nuevo.  El valor pedagógico está en la
\textbf{retroalimentación inmediata}: el alumno recibe el examen mientras
el material aún está fresco y puede cerrar sus huecos antes de que el curso
avance.

\begin{notaprofesor}[Por qué esta sesión importa]
  En cursos de ciencias exactas, los errores en los fundamentos se acumulan.
  Un alumno que no entiende el complemento a 2 hoy tendrá problemas con
  apuntadores, aritmética de memoria y tipos de datos más adelante.
  Esta pausa es la oportunidad de nivelar al grupo \textit{antes} de entrar
  a los temas de memoria y apuntadores (sesiones 11 en adelante).
\end{notaprofesor}

% =============================================================================
\section{Estructura sugerida de la sesión}
% =============================================================================

\begin{center}
\begin{tabular}{clc}
  \toprule
  \textbf{Tiempo} & \textbf{Actividad} & \textbf{Notas} \\
  \midrule
  10 min & Entregar exámenes y dar tiempo de lectura personal
         & Silencio; que cada quien lea su resultado \\
  50 min & Revisión pregunta por pregunta en el pizarrón
         & Ver sección~\ref{sec:revision} \\
  10 min & Preguntas abiertas y cierre
         & Aclarar dudas individuales \\
  \midrule
  10 min & (Opcional) Anticipo de la sesión 11
         & Solo si el tiempo lo permite \\
  \bottomrule
\end{tabular}
\end{center}

% =============================================================================
\section{Cómo conducir la revisión}
\label{sec:revision}
% =============================================================================

\begin{notaprofesor}[Orden recomendado]
  Revisa las preguntas de \textbf{menor a mayor dificultad}, no en el
  orden del examen.  Esto permite que los alumnos recuperen confianza
  con las correcciones fáciles antes de enfrentar las que más dolieron.

  Para cada pregunta:
  \begin{enumerate}
    \item Muestra la solución completa en el pizarrón, con el proceso.
    \item Pregunta al grupo: ``¿alguien lo resolvió diferente y llegó al
          mismo resultado?''  Valida métodos alternativos correctos.
    \item Señala el \textbf{error más frecuente} que viste al calificar
          y explica exactamente dónde está la confusión.
    \item No dediques más de 5--7 minutos por pregunta.
          Si una pregunta genera mucha discusión, agenda una sesión de
          asesoría fuera de clase.
  \end{enumerate}
\end{notaprofesor}

\begin{notaprofesor}[Errores que suelen generar más preguntas]
  Con base en el diseño del examen de la sesión 09, los temas que
  históricamente producen más dudas en la revisión son:
  \begin{itemize}
    \item \textbf{Complemento a 2}: confusión entre interpretar un patrón
          con signo vs.\ sin signo.
    \item \textbf{Fracciones binarias}: no saber determinar si una fracción
          es exacta sin hacer la conversión.
    \item \textbf{NOT bit a bit}: confundir \texttt{\textasciitilde x}
          con \texttt{-x} en lugar de \texttt{-(x+1)}.
    \item \textbf{Conversiones cruzadas} (e.g., hex $\to$ octal): no pasar
          por binario como puente y tratar de hacer la conversión directa.
  \end{itemize}
\end{notaprofesor}

% =============================================================================
\section{Después de la revisión}
% =============================================================================

\begin{notaprofesor}[Seguimiento]
  \begin{itemize}
    \item Si más de un tercio del grupo reprobó algún bloque temático,
          considera agregar un ejercicio corto de repaso al inicio de la
          sesión 11 antes de avanzar al nuevo tema.
    \item Los alumnos que quieran discutir su calificación individualmente
          deben hacerlo en asesoría, no durante la clase.
          Establece ese límite al inicio de la sesión.
  \end{itemize}
\end{notaprofesor}

\end{document}
